\documentclass{article}
\usepackage{graphicx} % Required for inserting images
\usepackage{amsmath}
\usepackage{amsthm}
\usepackage{amssymb}
\usepackage{mathrsfs}
\usepackage{enumitem}
\newtheorem{theorem}{Theorem}
\newtheorem{lemma}{Lemma}
\newtheorem{corollary}{Corollary}
\theoremstyle{remark}\newtheorem{remark}{Remark}

% Theorem with manually specified number, to match Goldrei's numbering
\theoremstyle{plain}
\newtheorem{bookthminner}{Theorem}
\newenvironment{bookthm}[1]
  {\renewcommand\thebookthminner{#1}\bookthminner}
  {\endbookthminner}

\newcommand{\T}[0]{\mathcal{T}}
\newcommand{\B}[0]{\mathcal{B}}
\newcommand{\C}[0]{\mathcal{C}}
\newcommand{\R}[0]{\mathbb{R}}
\newcommand{\Q}[0]{\mathbb{Q}}
\newcommand{\Z}[0]{\mathbb{Z}}
\newcommand{\N}[0]{\mathbb{N}}
\newcommand{\e}[0]{\epsilon}

\begin{document}

\noindent 2.38 Suppose that the formula $\phi$ is constructed by taking subformulas $\theta_1, \theta_2, ... , \theta_n$ in that order and joining them together only using the connective $\land$, with brackets inserted in such a way as to make $\phi$ a formula. Show that $\phi$ is true precisely for those truth assignments which make each of the subformulas $\theta_1, \theta_2, ... , \theta_n$ true. Deduce that any two such formulas (using the same $\theta_1, \theta_2, ... , \theta_n$ in that order) are logically equivalent. [\textit{Hint:} Use the version of mathematical induction with the hypothesis that the result holds for all $k \leq n$ where $n \geq 1$.]

% I need to show;
% 1. what we're inducting on
% 2. how \phi is built
% 3. what P(n) is (where P(n) is some property and n is some number)

\begin{proof}
    We proceed by induction on the number of $\theta$'s available in $\phi$.
    $\phi$ is constructed by taking subformulas $\theta_1, \theta_2, ... , \theta_n$ in that order and joining them together only using the connective $\land$, with brackets inserted in such a way as to make $\phi$ a formula.
    We will show $\phi$ is true precisely for those truth assignments which make each of the subformulas $\theta_1, \theta_2, ... , \theta_n$ true.
    
    \underline{Base Case:}
    Suppose that $\phi = \theta_1$. Let $v$ be any truth assignment. Therefore $v(\phi) = T \iff v(\theta_1) = T$.

    \underline{Inductive Hypothesis:}
    Let $n \geq 1$. Suppose that for every $k \leq n$ the following holds: if $\psi$ is a formula built from the formulas $\psi_1, \psi_2, \ldots, \psi_k$ in that order, joined only by the connective $\land$ with brackets inserted in such a way as to make $\psi$ a formula, then for every truth assignment $v$,
    \[
        v(\psi) = T \iff v(\psi_j) = T \text{ for all } j \text{ with } 1 \leq j \leq k.
    \]

    \underline{Inductive Step:}
    Let $\phi$ be built from the formulas $\theta_1, \theta_2, ..., \theta_n, \theta_{n+1}$ in that order, joined only by the connective $\land$ with brackets inserted in such a way as to make $\phi$ a formula.

     $\phi$ has a principal connective, so let $\theta_i$ be immediately to the left of the principal connective and let $\theta_{i+1}$ be immediately to the right of the principal connective where $1 \leq i \leq n$.
     
     Therefore $\phi$ can be rewritten as $\phi = \alpha \land \beta$ where $\alpha$ is built from the formulas $\theta_1, \theta_2, ... , \theta_i$ and $\beta$ is built from the formulas $\theta_{i + 1}, ... , \theta_n, \theta_{n + 1}$ in that order, joined only by the connective $\land$ with brackets inserted in such a way as to make them each valid subformulas of $\phi$.yea

     Since $i \leq n$ and $(n + 1) - i \leq n$ the inductive hypothesis applies to $\alpha, \beta$.

     Let $v$ be any truth assignment. Since $\phi = \alpha \land \beta$, the truth table for $\land$ gives
     \[
         v(\phi) = T \iff v(\alpha) = T = v(\beta).
     \]
     By the inductive hypothesis with $k = i$ and $\psi_1, \ldots, \psi_i := \theta_1, \ldots, \theta_i$,
     \[
         v(\alpha) = T \iff v(\theta_j) = T \text{ for all } j \text{ with } 1 \leq j \leq i.
     \]
     By the inductive hypothesis with $k = (n+1) - i$ and $\psi_1, \ldots, \psi_{(n+1)-i} := \theta_{i+1}, \ldots, \theta_{n+1}$,
     \[
         v(\beta) = T \iff v(\theta_j) = T \text{ for all } j \text{ with } i+1 \leq j \leq n+1.
     \]
     Combining these three biconditionals,
     \[
         v(\phi) = T \iff v(\theta_j) = T \text{ for all } j \text{ with } 1 \leq j \leq n+1,
     \]
     which is precisely the claim for $n+1$. This completes the induction.
\end{proof}

\noindent\underline{Deduction:}
Suppose $\phi$ and $\phi'$ are both built from the formulas $\theta_1, \theta_2, \ldots, \theta_n$ in that order, joined only by the connective $\land$ with brackets inserted in such a way as to make each a formula. Let $v$ be any truth assignment. By the result above applied to each formula,
\[
    v(\phi) = T \iff v(\theta_j) = T \text{ for all } j \text{ with } 1 \leq j \leq n \iff v(\phi') = T.
\]
If $v(\phi) = T$, then $v(\phi') = T$. If $v(\phi) = F$, then $v(\phi) \neq T$, so $v(\phi') \neq T$, and by bivalence $v(\phi') = F$. In either case $v(\phi) = v(\phi')$. Since $v$ was arbitrary, $\phi \equiv \phi'$.

\end{document}

